\originalchapter{三能级相干与集体辐射}
\originalsection{Raman耦合与暗态}
取$\Lambda$系统两个稳定态$\ket1,\ket2$和激发态$\ket3$. 在双光子共振的旋转系中,
\[\frac H\hbar=\begin{pmatrix}0&0&\Omega_p^*/2\\0&0&\Omega_c^*/2\\\Omega_p/2&\Omega_c/2&-\Delta\end{pmatrix}.\]
大单光子失谐$|\Delta|\gg|\Omega_p|,|\Omega_c|,\Gamma$时, 令$\dot c_3\simeq0$, 得$c_3=(\Omega_pc_1+\Omega_cc_2)/(2\Delta)$. 代回前两式得有效二能级哈密顿量, 对角AC位移为$|\Omega_{p,c}|^2/(4\Delta)$, 非对角为$\Omega_p^*\Omega_c/(4\Delta)$. 所以有效Rabi频率$|\Omega_p\Omega_c|/(2|\Delta|)$, 散射误差则按$\Gamma\Omega^2/\Delta^2$缩小. 增大失谐若同时提高功率保持Raman速率, 需重新评估功率与其他能级的作用.

精确暗态为$\ket D=(\Omega_c\ket1-\Omega_p\ket2)/\sqrt{|\Omega_c|^2+|\Omega_p|^2}$, 因激发态耦合振幅相消, $H\ket D=0$. 先开控制场、后开泵浦场, 再逐渐关闭控制场, $\ket D$由$\ket1$转到$-\ket2$, 是STIRAP的反直觉脉冲顺序. 绝热条件由暗态与亮态的最小能隙决定, 且双光子失谐必须足够小. 有限退相干会破坏暗态, 所以“无激发态布居”是理想模型的结论, 不是对任何脉冲都成立.

\originalsection{EIT与色散}
弱探针使大部分原子留在$1$. 令光学相干衰减$\gamma_{31}$、基态相干衰减$\gamma_{21}$、探针与双光子失谐$\Delta_p,\delta$, 线性稳态方程为
\[0=-(\gamma_{31}-\ii\Delta_p)\rho_{31}+\frac\ii2\Omega_p+\frac\ii2\Omega_c\rho_{21},\quad0=-(\gamma_{21}-\ii\delta)\rho_{21}+\frac\ii2\Omega_c^*\rho_{31}.\]
消去$\rho_{21}$得
\[\rho_{31}=\frac{\ii\Omega_p(\gamma_{21}-\ii\delta)/2}{(\gamma_{31}-\ii\Delta_p)(\gamma_{21}-\ii\delta)+|\Omega_c|^2/4}.\]
$\gamma_{21}\to0,\delta=0$时分子为零, 吸收消失. 附近的实部仍有陡色散, 导致群速度改变. 透明窗口来自量子振幅干涉, 不等于减少所有频率的原子吸收; 有限窗口、脉冲带宽和基态退相干限制可储存或延迟的波包.

\originalsection{Dicke态与超辐射}
若$N$个相同原子位于远小于波长的体积中, 场不能分辨发射者, 跳跃算符为$J_- =\sum_i\sigma_-^{(i)}$. 对称角动量态$\ket{J=N/2,M}$满足
\[J_-\ket{J,M}=\sqrt{(J+M)(J-M+1)}\ket{J,M-1}.\]
所以辐射率为$\Gamma(J+M)(J-M+1)$. 全激发态起始率$N\Gamma$, 在激发约一半时峰值约$N^2\Gamma/4$. 这说明强度随$N^2$增长可由集体量子相关产生, 不必先假定每个原子都有非零偶极平均.

两个原子的单激发对称态$(\ket{eg}+\ket{ge})/\sqrt2$发射率$2\Gamma$, 反对称态在完全相同耦合的理想模型中为暗态. 大样品存在传播相位、几何方向性与再吸收, 不能直接套小样品总率. 集体辐射、经典相干散射和受激发射的实验表现可能相似, 其初态、场占据和相关函数却不同.

\originalsection{扩展样品的相位与有效耦合}
8.421第25讲第20--22页保留传播相位来讨论大样品. 对位置$\mathbf r_j$的原子, 方向$\mathbf k$的发射算符是$J_-(\mathbf k)=\sum_j\e^{-\ii\mathbf k\cdot\mathbf r_j}\sigma_-^{(j)}$. 取带空间相位的单激发态
\[\ket{1_{\mathbf q}}=\frac1{\sqrt N}\sum_j
 \e^{\ii\mathbf q\cdot\mathbf r_j}\ket{g\cdots e_j\cdots g},\qquad
 \frac{\dd\Gamma_{\mathbf q}}{\dd\Omega}
 =\frac{\dd\Gamma_1}{\dd\Omega}\frac1N
 \left|\sum_j\e^{\ii(\mathbf q-\mathbf k)\cdot\mathbf r_j}\right|^2.\]
在$\mathbf k=\mathbf q$时振幅为$\sqrt N$, 该方向微分率增$N$倍; 偏离方向的路径相位逐渐相消. 若位置独立取自归一空间分布, 结构因子平均为$1+(N-1)|F(\mathbf q-\mathbf k)|^2$, $F(\mathbf p)=\langle\e^{\ii\mathbf p\cdot\mathbf r}\rangle$. 因而增强并非遍及整个$4\pi$. 例如均匀半径$R$球的$F=3[\sin u-u\cos u]/u^3$, $u=|\mathbf p|R$, 主瓣角宽约$1/(kR)$, 对应固角量级$\lambda^2/R^2$. 这是几何如何进入总率的一个可核算模型.

原课的长样品用相位匹配轴向模式和有效横截面$A_{\rm eff}$估计相干固角$\Delta\Omega\sim\lambda^2/A_{\rm eff}$. 把偶极方向权重和模式重叠纳入$\beta$（各向同性粗估为$\Delta\Omega/4\pi$）, 定义单原子在该共同辐射通道中的率$\Gamma_c=\beta\Gamma$. 只保留一个相位匹配通道及其余近独立通道的Markov有效模型为
\[\dot\rho=\Gamma_c\mathcal D[J_-(\mathbf q)]\rho
 +\Gamma_{\rm bg}\sum_j\mathcal D[\sigma_-^{(j)}]\rho,\qquad
 \mathcal D[L]\rho=L\rho L^\dagger-\tfrac12\{L^\dagger L,\rho\},\]
其中$\Gamma_{\rm bg}\simeq(1-\beta)\Gamma$. 给每个激发基矢乘$\e^{\ii\mathbf q\cdot\mathbf r_j}$只是局域相位旋转, 所以带相位的$m$激发Dicke态仍有$\langle J_+J_-\rangle=m(N-m+1)$. 共同通道率为$\Gamma_c m(N-m+1)$, 独立背景率为$\Gamma_{\rm bg}m$. 对$m=1$, 总率约$\Gamma[1+\beta(N-1)]$; 半激发峰值在共同通道约$\beta\Gamma N^2/4$. 小样品所有通道同相时须改为同一个集体跳跃, 才恢复原来的全空间$N^2$峰值.

若$N=nA_{\rm eff}L$, 轴向增强量$N\beta$按$n\lambda^2L$增长, 与共振光学厚度$b_0=n\sigma_0L$同尺度（闭合偶极跃迁$\sigma_0=3\lambda^2/(2\pi)$）. 这解释第25讲为何把增强联系到沿程增益. 这里的$\beta$必须由实际结构因子、纵向相位匹配和模式几何确定, 不是任意长度都用同一孔径常数. 非轴向发射的相位不同, 会把态带出该Dicke梯级; 再吸收、多个横模、失谐和Doppler去相干也可破坏此简单模型. 传播时间$L/c$不再短于集体衰减时间时还需含延迟的场方程, 不可把这个单通道Markov式当完整Maxwell--Bloch解.
